% =====================================================================
% Instrukcja laboratoryjna — XPlanar / TwinCAT 3
% Szkielet dokumentu: rozdziały = zagadnienia laboratorium, placeholdery.
% Styl wizualny zbliżony do dokumentacji Beckhoff (InfoSys).
% =====================================================================
\documentclass[11pt,a4paper]{report}

\usepackage[utf8]{inputenc}
\usepackage[T1]{fontenc}
\usepackage[polish]{babel}
\usepackage[scaled=0.92]{helvet}
\renewcommand{\familydefault}{\sfdefault}
\usepackage[margin=2.4cm,headheight=14pt]{geometry}
\usepackage{microtype}
\usepackage{graphicx}
\usepackage{float}
\usepackage{xcolor}
\usepackage{fancyhdr}
\usepackage{titlesec}
\usepackage{enumitem}
\usepackage{booktabs}
\usepackage[most]{tcolorbox}
\usepackage[colorlinks=true,linkcolor=bkorange,urlcolor=bkorange,citecolor=bkorange]{hyperref}

\graphicspath{{zdjęcia/}}

% --- Paleta (Beckhoff: pomarańcz + szarości) ---
\definecolor{bkorange}{RGB}{232,90,10}
\definecolor{bkdark}{RGB}{51,51,51}
\definecolor{bkgray}{RGB}{240,240,240}
\definecolor{bkblue}{RGB}{0,106,163}
\definecolor{bknote}{RGB}{230,242,250}
\definecolor{bkwarn}{RGB}{255,244,214}
\definecolor{bkwarnbd}{RGB}{214,158,0}

\AtBeginDocument{\color{bkdark}}
\setlength{\parindent}{0pt}
\setlength{\parskip}{5pt}
\setlist{itemsep=1pt,topsep=3pt}
\emergencystretch=3em

% --- Nagłówki ---
\titleformat{\chapter}[block]
  {\normalfont\bfseries\color{bkorange}}
  {\fontsize{34}{36}\selectfont\thechapter}{0.5em}
  {\Huge\color{bkdark}}
  [\vspace{2pt}{\color{bkorange}\titlerule[2pt]}]
\titlespacing*{\chapter}{0pt}{-10pt}{22pt}
\titleformat{\section}{\Large\bfseries\color{bkorange}}{\thesection}{0.6em}{}
\titlespacing*{\section}{0pt}{16pt}{6pt}
\titleformat{\subsection}{\large\bfseries\color{bkdark}}{\thesubsection}{0.6em}{}
\titleformat{\subsubsection}{\normalsize\bfseries\color{bkorange}}{}{0em}{}
\titlespacing*{\subsubsection}{0pt}{12pt}{4pt}
\setcounter{tocdepth}{3}
\newcommand{\podrozdzial}[1]{\subsection*{#1}\addcontentsline{toc}{subsection}{#1}}
\newcommand{\podpodrozdzial}[1]{\subsubsection*{#1}\addcontentsline{toc}{subsubsection}{#1}}

% --- Nagłówek i stopka ---
\pagestyle{fancy}
\fancyhf{}
\renewcommand{\headrulewidth}{0.8pt}
\renewcommand{\headrule}{\hbox to\headwidth{\color{bkorange}\leaders\hrule height \headrulewidth\hfill}}
\fancyhead[L]{\small\color{bkdark}XPlanar --- TwinCAT 3}
\fancyhead[R]{\small\color{bkdark}\leftmark}
\fancyfoot[L]{\small\color{gray}Instrukcja laboratoryjna}
\fancyfoot[R]{\small\color{bkorange}\textbf{\thepage}}
\renewcommand{\chaptermark}[1]{\markboth{#1}{}}
\fancypagestyle{plain}{\fancyhf{}\fancyfoot[R]{\small\color{bkorange}\textbf{\thepage}}\renewcommand{\headrulewidth}{0pt}}

% --- Ramki w stylu InfoSys: kolorowy pasek z lewej ---
\tcbset{beckbox/.style={enhanced,breakable,frame hidden,boxrule=0pt,sharp corners,
  borderline west={3pt}{0pt}{bkorange},left=10pt,right=8pt,top=6pt,bottom=6pt}}

\newcommand{\opisDoUzupelnienia}[1]{%
  \begin{tcolorbox}[beckbox,colback=bkgray,borderline west={3pt}{0pt}{bkdark!50}]
  \textbf{[DO UZUPEŁNIENIA]} Opis czynności oraz nazewnictwo bloków, metod i parametrów
  zgodne z dokumentacją Beckhoff (\texttt{Tc3\_XPlanarStandard}, \texttt{Tc3\_XPlanarUtility}). #1
  \end{tcolorbox}}

\newcommand{\uwaga}[1]{%
  \begin{tcolorbox}[beckbox,colback=bknote,borderline west={3pt}{0pt}{bkblue}]
  \textbf{\color{bkblue}Note}\par #1
  \end{tcolorbox}}

\newcommand{\zadanie}[2]{%
  \begin{tcolorbox}[beckbox,colback=bkwarn,borderline west={3pt}{0pt}{bkwarnbd}]
  \textbf{\color{bkorange}#1}\par #2
  \end{tcolorbox}}

% --- Placeholder na zrzut ekranu ---
\newcommand{\zrzutEkranu}[3][0.8\linewidth]{%
  \begin{figure}[H]
    \centering
    \IfFileExists{zdjęcia/#2}{%
      \includegraphics[width=#1,height=0.42\textheight,keepaspectratio]{#2}%
    }{%
      \begin{tcolorbox}[enhanced,width=#1,height=5.5cm,colback=bkgray,colframe=bkdark!30,
        boxrule=0.6pt,sharp corners,halign=center,valign=center]
        \itshape Miejsce na zrzut ekranu\\[0.4em]\normalfont\ttfamily\small zdjęcia/#2
      \end{tcolorbox}%
    }
    \caption{#3}
  \end{figure}}

\title{\textbf{Instrukcja laboratoryjna}\\\vspace{0.4em}\large XPlanar --- TwinCAT 3 (IEC~61131-3)}
\author{}
\date{}

\begin{document}

\maketitle
\tableofcontents
\newpage

\chapter{Wstęp}
\textbf{Forma:} prowadzący pokazuje technologię wizualnie (na maszynie), nie na slajdach.

\subsection*{Komentarz / zalecenia}
\begin{itemize}
  \item Student musi otrzymać wiedzę przed przystąpieniem do laboratorium numer 1.
  \item Wytłumaczenie musi wystąpić \textbf{wizualnie} --- ilość wiedzy jest zbyt duża, by przedstawić ją na wykładzie, którego slajdy nie odzwierciedlają ani nie uczą realnej pracy przy maszynie.
  \item Z doświadczenia: po zajęciach z PLC lub jakichkolwiek zajęciach z automatyki przemysłowej studenci nie dochodzą do wyjaśnienia zasady działania ani postępowania z maszyną. Typowy student nawet nie ma pojęcia, czym \textit{dokładnie} jest PLC.
  \item Wizualne przekazanie wiedzy, co można zrobić z danym sprzętem --- choćby na podstawie PLC nie od Beckhoffa, a np. Delty/Siemensa, na realnych przykładach.
  \item To kluczowe do efektywnego przeprowadzenia laboratoriów, ponieważ systemy Beckhoffa cechują się wysokim poziomem wejścia relatywnie do innych środowisk, jak TIA Portal czy ISPSoft.
  \item Podkreślić, że na rok 2026 platforma XPlanar Beckhoff jest technologią jeszcze nie w 100\% dopracowaną: na razie jest rozwijana i mało jest systemów zaimplementowanych, jednakże jest to rozwiązanie bardzo rzeczowe / obiecujące \textit{(do doprecyzowania sformułowanie)}.
\end{itemize}

\section{Jak działa XPlanar i jego możliwości}
\begin{itemize}
  \item jak działa,
  \item jakie są maksymalne zakresy działania,
  \item jakie typy i rozmiary,
  \item możliwości i miejsca zastosowania,
  \item ograniczenia moverów.
\end{itemize}
\opisDoUzupelnienia{}

\section{Szafa sterownicza}
\zrzutEkranu[0.8\linewidth]{01-wstep/szafa-sterownicza/szafa-sterownicza-widok-ogolny.png}{Szafa sterownicza --- widok ogólny (zdjęcie do dodania)}
\begin{itemize}
  \item co znajduje się w szafie,
  \item jak elementy szafy komunikują się ze sobą i z peryferiami --- EtherCAT.
\end{itemize}
\opisDoUzupelnienia{}

\section{Środowisko TwinCAT XAE Shell}
Projekty TwinCAT 3 tworzymy w środowisku \textbf{TcXaeShell 64-bit} (wersja 7.21.11, oparte na Visual Studio 2022) wraz z \textbf{TwinCAT XAE 4026.27}.

\podrozdzial{Okno główne}
Okno główne składa się z paska menu i paska narzędzi, Solution Explorera z drzewem projektu, obszaru edytora, w którym otwieramy wybrane elementy projektu, oraz okien pomocniczych (m.in.\ Error List i Output). 
\podrozdzial{Pasek menu i narzędzi}
\zrzutEkranu[0.95\linewidth]{01-wstep/srodowisko-tcxaeshell/tcxaeshell-pasek-narzedzi.png}{Edytor TcXaeShell --- pasek menu i pasek narzędzi}

Z paska, z którego będziemy realnie korzystać, od lewej:
\begin{enumerate}
  \item \textbf{wybór wersji TwinCAT}, przez którą się komunikujemy, a następnie wersja aktywna (\textit{Active version}),
  \item \textbf{Activate Configuration} --- wgranie konfiguracji projektu na wybrany target,
  \item \textbf{Run Mode} i \textbf{Config Mode} --- przełączanie trybu pracy runtime'u (zielona --- Run, niebieska --- Config),
  \item \textbf{aktywny projekt} (\texttt{StarterkitBase}) --- lista rozwijana projektów w rozwiązaniu,
  \item \textbf{target} (\texttt{<Local>}) --- system, na którym uruchamiamy projekt.
\end{enumerate}



Sterownik zewnętrzny wybieramy po nazwie lub numerze seryjnym; wymaga to wcześniej dodanej trasy komunikacji (\textit{Edit Routes w Zasobniku}).

\podrozdzial{Solution Explorer}
Solution Explorer dzieli projekt na części: SYSTEM, MOTION, PLC, SAFETY, C++, VISION, ANALYTICS i I/O. W laboratorium korzystamy z SYSTEM, MOTION, I/O i PLC.

\podpodrozdzial{SYSTEM}
SYSTEM to konfiguracja komputera sterującego, na którym uruchamiamy projekt.

\zrzutEkranu[0.4\linewidth]{01-wstep/czesci-projektu/solution-explorer-system.png}{Solution Explorer --- SYSTEM}

\textbf{License} --- licencje TwinCAT.

\textbf{Real-Time} --- w tej zakładce definiujemy rdzenie procesora, które pracują, aby system działał szybko.

\zrzutEkranu[0.95\linewidth]{01-wstep/czesci-projektu/solution-explorer-system-realtime.png}{Solution Explorer --- SYSTEM, Real-Time (Settings)}

Rdzenie \textit{shared} (współdzielone) TwinCAT dzieli z Windows, a \textit{Core Limit} określa, jaką część czasu rdzenia może zająć TwinCAT. Rdzenie \textit{isolated} (izolowane) są przydzielone wyłącznie TwinCAT i Windows z nich nie korzysta. Po zmianie parametrów trzeba zapisać konfigurację (\textit{Set on Target}).

\textbf{TcCOM Objects} --- obiekty, w których żyje system XPlanar. Obiekt \texttt{XPlanar} zawiera:
\begin{itemize}
  \item \texttt{Part1} --- kafelki, po których poruszają się movery,
  \item \texttt{Area1} --- obszar ruchu, czyli zdefiniowane środowisko, w którym movery mogą się poruszać,
  \item movery, nazwane swoim numerem \textbf{BTN} (np.\ \texttt{0000kyte}).
\end{itemize}
BTN (\textit{Beckhoff Traceability Number}) to numer fabryczny movera z jego naklejki. Mover jest tylko magnesem, więc identyfikuje go zderzak (czarna ramka); dzięki BTN system wie, który mover jest którym obiektem w projekcie.

\podpodrozdzial{MOTION}
\texttt{XPlanar MC Project} zawiera obiekty ruchu, które są pośrednikiem między programem PLC a światem rzeczywistym.

\zrzutEkranu[0.5\linewidth]{01-wstep/czesci-projektu/solution-explorer-motion.png}{Solution Explorer --- MOTION}

W jego skład wchodzą:
\begin{itemize}
  \item \texttt{Axes} --- po jednej osi na każdy mover (to nie są osie silników). Oś jest zlinkowana z obiektem movera w TcCOM oraz ze zmiennymi PLC (\texttt{PlcToMc}: komendy, \texttt{McToPlc}: stany i pozycja); linki PLC wykonujemy sami (\textit{Change Axis PLC Links}),
  \item \texttt{Groups} --- grupa \textbf{XPlanar} (sprawdzanie kolizji movery--tracki--środowisko), \textbf{Environment} (powierzchnia kafelków, czyli obszar, w którym się poruszamy) oraz \textbf{tracki} (\textit{Track1--3}, „ulice'' dla moverów).
\end{itemize}
Tylko osie moverów mają odpowiednik w sprzęcie. Grupa, Environment i tracki są czysto programowe, więc linkuje się je tylko z PLC (\texttt{MCTOPLC\_STD}), aby PLC mogło nimi sterować.

\podpodrozdzial{I/O}
I/O opisuje fizyczne urządzenia i powiązania między nimi a resztą projektu.

\zrzutEkranu[0.4\linewidth]{01-wstep/czesci-projektu/solution-explorer-io.png}{Solution Explorer --- I/O}
\zrzutEkranu[0.6\linewidth]{01-wstep/czesci-projektu/solution-explorer-io-mappings.png}{Solution Explorer --- I/O, Mappings}

Na zrzutach widać:
\begin{itemize}
  \item trzy urządzenia EtherCAT: Device 1, Device 2 i Device 3,
  \item w Device 1 sprzęgacz EK1200,
  \item w Device 2 terminale APS4322-7000 o numerach 5, 7, 9, 11, 13 i 15, czyli kafelki XPlanara,
  \item w Device 3 terminale APS4322-7000 o numerach 17, 19, 21, 23, 25 i 27, również kafelki,
  \item w każdym urządzeniu EtherCAT elementy wspólne: Image, Image-Info, SyncUnits, Inputs, Outputs i InfoData,
  \item urządzenie USB: Device 6,
  \item Mappings, czyli wszystkie linki w jednym miejscu: obiekty moverów 0000kyte, 0000kyto, 0000kytm i 0000kytn powiązane ze swoimi osiami, a także Context, XPlanar, Environment i Track1--3 powiązane z instancją projektu PLC.
\end{itemize}

\podrozdzial{Ikona TwinCAT w zasobniku systemu}
Menu ikony TwinCAT w zasobniku systemu pozwala m.in. uruchomić środowisko (TcXaeShell 32/64-bit, Visual Studio), uruchomić lub zrestartować runtime (\texttt{Start/Restart}), przełączyć go w tryb Config (\texttt{Config}) oraz otworzyć ustawienia routera ADS (\texttt{Router}: \texttt{Info}, \texttt{Cleanup}, \texttt{Change AMS NetId}, \texttt{Edit Routes}). Stan runtime'u widać po kolorze ikony (niebieska --- Config, zielona --- Run); szczegóły: \href{https://infosys.beckhoff.com/content/1033/tc3_userinterface/2867291659.html}{Information bar} w Infosys.

\zrzutEkranu[0.4\linewidth]{01-wstep/srodowisko-tcxaeshell/tcxaeshell-zasobnik-systemu.png}{Menu ikony TwinCAT w zasobniku systemu}

\section{Kod demo}
\zrzutEkranu[0.4\linewidth]{01-wstep/kod-demo-sk12/solution-explorer-plc.png}{Solution Explorer --- PLC}

W tym projekcie skupiamy się tylko na folderze \texttt{2\_Application}.

\podrozdzial{Podstawy programowania PLC}
Językiem programowania jest \textbf{ST} (\textit{Structured Text}).

\textbf{Zasady.} Ze względu na to, że jest to PLC, nasz program wykonuje się \textbf{cyklicznie}, podobnie jak w ladderze czy w języku blokowym: zadanie \texttt{PlcTask} wywołuje program główny w każdym cyklu (tu co 10 ms). Dlatego kod nie może czekać w pętli; długie czynności rozkładamy na kroki maszyny stanów (\texttt{CASE}), wykonywane w kolejnych cyklach. Poza tym:
\begin{itemize}
  \item zmienną deklarujemy przed użyciem, w bloku \texttt{VAR \ldots END\_VAR}, z typem i opcjonalną wartością początkową,
  \item instrukcję kończymy średnikiem; przypisanie to \texttt{:=}, porównanie to \texttt{=}.
\end{itemize}

\textbf{Jednostki organizacyjne (POU).} Program składa się z jednostek organizacyjnych (\textit{Program Organization Unit}), które dodajemy w Solution Explorer:
\begin{itemize}
  \item \textbf{Program} --- jednostka wywoływana bezpośrednio przez zadanie; jest jedna i zachowuje stan między cyklami,
  \item \textbf{Function Block} --- jednostka z własną pamięcią: tworzymy jej instancje, a każda pamięta swój stan. Może mieć metody i właściwości, dziedziczyć po innym bloku i implementować interfejsy,
  \item \textbf{Function}  --- jednostka bez pamięci, która dla tych samych argumentów zwraca ten sam wynik,
  \item \textbf{Method}  --- funkcja należąca do bloku funkcyjnego lub programu, wywoływana na jego instancji, z dostępem do jego zmiennych,
  \item \textbf{Property} (właściwość) --- z zewnątrz działa jak zmienna bloku, ale odczyt i zapis wykonuje kod (\textit{Get} i \textit{Set}),
  \item \textbf{Action}  --- dodatkowy fragment kodu bloku lub programu, bez własnych argumentów, pracujący na jego zmiennych,
  \item \textbf{Interface}  --- lista metod i właściwości bez kodu; blok, który go implementuje, musi je zawierać, dzięki czemu różne bloki można traktować jednakowo,
  \item \textbf{DUT} (\textit{Data Unit Type}) --- definicja własnego typu danych: struktury, wyliczenia lub aliasu,
  \item \textbf{GVL} (\textit{Global Variable List}) --- lista zmiennych i stałych globalnych, dostępnych w całym projekcie.
\end{itemize}

\textbf{Dobre praktyki nazewnictwa.} Nazwy zmiennych mają prefiks oznaczający typ: \texttt{b} --- \texttt{BOOL}, \texttt{n} --- liczba całkowita, \texttt{f} --- zmiennoprzecinkowa, \texttt{s} --- \texttt{STRING}, \texttt{fb} --- instancja bloku funkcyjnego, \texttt{ip} --- interfejs. Typy i listy mają prefiks rodzaju: \texttt{FB\_}, \texttt{ST\_} (struktura), \texttt{E\_} (wyliczenie), \texttt{I\_} (interfejs), \texttt{GVL\_}; stałe globalne zaczynamy od \texttt{gc\_}.

\podrozdzial{Program}
\zrzutEkranu[0.4\linewidth]{01-wstep/kod-demo-sk12/solution-explorer-plc-kod.png}{Solution Explorer --- PLC, kod aplikacji}

\begin{itemize}
  \item \texttt{Machine} --- \texttt{FB\_MotionSystem} to główna maszyna stanów systemu: zatrzymanie, reset, przygotowanie i praca. W przygotowaniu sprawdza zasilanie kafelków, wykrywa i identyfikuje movery, określa ich orientację, włącza je i wprowadza na tory. \texttt{FB\_MoverTarget} wysyła movery do zadanych pozycji,
  \item \texttt{Mover} --- \texttt{FB\_ApplicationMover} to obiekt aplikacyjny jednego movera: w każdym cyklu wywołuje standardowy obiekt movera i udostępnia jego pozycję wizualizacji,
  \item \texttt{Station} --- \texttt{FB\_StationBase} to wspólna baza dla pięciu stacji. Każda stacja jest maszyną stanów z metodami \texttt{Init}, \texttt{Clear} i \texttt{Process}: czeka na movera, wysyła go do pozycji obróbki, wykonuje swój proces, wybiera następną stację i odsyła movera po torze. Nazwy stacji 4 i 5 (\texttt{SW\_Separate}, \texttt{SW\_Merge}) wskazują na rozdzielanie i łączenie torów,
  \item \texttt{Tracks} --- \texttt{FB\_TrackBase} i trzy segmenty torów, które po nim dziedziczą; metoda \texttt{Initialize} buduje tablicę punktów toru,
  \item \texttt{GVL\_Application} --- stałe aplikacji: 4 movery, 3 tory i 5 stacji.
\end{itemize}

\podrozdzial{References}
\zrzutEkranu[0.4\linewidth]{01-wstep/kod-demo-sk12/solution-explorer-plc-references.png}{Solution Explorer --- PLC, References}

References to biblioteki dołączone do projektu: gotowe bloki i typy, z których korzysta program. Używamy między innymi \texttt{Tc3\_XPlanarStandard}, która nie jest ogólnie udostępniana przez Beckhoff. Dzięki niej pisanie programu jest wyższego poziomu: zamiast niskopoziomowych komend ruchu używamy gotowych obiektów systemu i movera, np.\ wywołania \texttt{std.cycle()}.

\podrozdzial{Linki Inputs i Outputs}
\zrzutEkranu[0.8\linewidth]{01-wstep/kod-demo-sk12/solution-explorer-plc-instance-linki.png}{Solution Explorer --- instancja PLC (Inputs, Outputs)}

Instancja projektu PLC ma dwie listy zmiennych wymienianych z resztą systemu w każdym cyklu zadania. \textbf{PlcTask Inputs} to dane wchodzące do PLC: przyciski pierścienia CX, a także statusy z MOTION (\texttt{MCTOPLC}) dla Environment, grupy, czterech moverów i trzech torów. \textbf{PlcTask Outputs} to dane wychodzące z PLC: komendy dla czterech moverów (\texttt{PLCTOMC}). \textbf{Link} to połączenie zmiennej PLC ze zmienną drugiej strony, np.\ osi w MOTION; dzięki niemu wartości są kopiowane automatycznie w każdym cyklu, a PLC może sterować moverami. Zmienne przyciskowe w programie są zadeklarowane jako wejścia (\texttt{AT \%I*}), czyli bez stałego adresu, który nadaje dopiero link.

\section{Odpalenie projektu}
\podrozdzial{Rodzaje pracy: Normal, BasicSimulation, AdvancedSimulation}
Tryb ustawiamy w obiekcie XPlanar: \texttt{SYSTEM} $\rightarrow$ \texttt{TcCOM Objects} $\rightarrow$ \texttt{XPlanar} $\rightarrow$ \texttt{Parameter (Init)} $\rightarrow$ \texttt{OperationMode}.

\zrzutEkranu[0.9\linewidth]{01-wstep/odpalenie-projektu/tccom-xplanar-parameter-init-operationmode.png}{Obiekt XPlanar --- Parameter (Init), wybór OperationMode}

\begin{center}\small
\begin{tabular}{@{}p{2.7cm}p{3.2cm}p{4.1cm}p{4.4cm}@{}}
\toprule
\textbf{Tryb} & \textbf{Sprzęt} & \textbf{Czas rzeczywisty} & \textbf{Obliczenia}\\
\midrule
\texttt{Normal} & wymagany & wymagany & rzeczywisty system\\
\midrule[0.3pt]
\texttt{BasicSimulation} & praca na komputerze & nie musi być spełniony, można symulwoać na jednym rdzeniu& symulowany ruch moverów, bez obliczania prądów i sprzężenia zwrotnego\\
\midrule[0.3pt]
\texttt{Advanced-\allowbreak Simulation} & ? & wymagany; układ zadań jak w rzeczywistej konfiguracji (czas cyklu, rozłożenie obciążenia na rdzenie) & aktywne obliczanie prądów i sprzężenia zwrotnego, bliżej trybu \texttt{Normal}\\
\bottomrule
\end{tabular}
\end{center}




\podrozdzial{Odpalenie w trybie Normal na PLC}
Tryb \texttt{Normal} to praca z rzeczywistym sprzętem na sterowniku XPlanara.
\begin{enumerate}
  \item Otwieramy rozwiązanie (aktualnie: \texttt{StarterkitBase}) w TcXaeShell.
  \item Wybieramy target: sterownik XPlanara (np.\ \texttt{CX-94E797}); wymaga to dodanej wcześniej trasy komunikacji.
  \item Sprawdzamy, czy \texttt{OperationMode} to \texttt{Normal}.
  \item Sprawdzamy, urządzenia I/O są enablowane.
  \item Upewniamy się, że movery stoją na kafelkach.
  \item Wykonujemy \textit{Activate Configuration}, logujemy się do PLC i uruchamiamy program.
  \item Maszynę obsługujemy z wizualizacji: Start, Stop, wprowadzenie w stan Execute.
\end{enumerate}

\podrozdzial{Odpalenie w symulacji}
W symulacji projekt działa bez maszyny: symulowany jest sprzęt XPlanar (kafelki i movery), a program PLC wykonuje się normalnie. Wspólne kroki:
\begin{enumerate}
  \item w obiekcie XPlanar ustawiamy \texttt{OperationMode} na \texttt{BasicSimulation} lub \texttt{AdvancedSimulation},
  \item urządzenia I/O ustawiamy na disable albo na simulation mode.
  \item wykonujemy \textit{Activate Configuration}, logujemy się do PLC i uruchamiamy program,
  \item zmiany w symulacji obserwujemy na dwa sposoby:
  \begin{itemize}
    \item \textit{TwinCAT} $\rightarrow$ \textit{MC3-XPlanar Environment View},
    \item \textit{TwinCAT} $\rightarrow$ \textit{XPlanar} $\rightarrow$ \textit{Configurator} $\rightarrow$ \textit{Live View}; po zmianach w projekcie klikamy \textit{Reload},
  \end{itemize}
\end{enumerate}

\podpodrozdzial{Symulacja na PLC Beckhoffa}
Jako target wybieramy sterownik. Runtime jest na nim już zainstalowany i skonfigurowany pod czas rzeczywisty, więc kroki dotyczące wirtualizacji i pakietów z następnego podrozdziału nie mają zastosowania. \texttt{BasicSimulation} nie wymaga czasu rzeczywistego, a \texttt{AdvancedSimulation} wymaga rzeczywistego układu zadań, który sterownik zapewnia. Przy niezgodności ustawień Real-Time wykonujemy \textit{Read from Target}.

\podpodrozdzial{Symulacja na zewnętrznym PC}
Jako target wybieramy \texttt{<Local>}. Komputer musi spełniać wymagania:
\begin{itemize}
  \item zainstalwoanwego środowiska TwinCAT w pdoanej konfigruacji. TODO DODOAĆ KONFIURACJE WYŁAUMCZYĆ PBORANIE NA KOMPUTER
  \item wyłączona wirtualizacja: Hyper-V, WSL, VBS (Memory integrity) i Windows Hello Enhanced Sign-in; zmiana wymaga uprawnień administratora i restartu,
  \item biblioteka \texttt{Tc3\_XPlanarStandard} w wersji, której wymaga projekt (References),
  \item ustawienia Real-Time dopasowane do komputera: \textit{Read from Target}, potem przypisanie zadań do istniejących rdzeni; dla \texttt{BasicSimulation} wystarczy jeden rdzeń, dla \texttt{AdvancedSimulation} układ zadań jak w rzeczywistej konfiguracji.
\end{itemize}


\section{Zadania}
\podrozdzial{Zadanie 1}
Uruchomienie projektu na PLC lub przez własny komputer w trybie Normal, odpalić z visualizacji ruch moverów 
\podrozdzial{Zadanie 2}
Odpalenie Symualcji w trybie BasicSimulation na PLC lub przez własny komputer, odpalić z visualizacji ruch moverów oraz sprawdic ruch

\subsection*{Dokumentacja Beckhoff (Infosys)}
\begin{itemize}
  \item \href{https://infosys.beckhoff.com/content/1033/tf5430_tc3_xplanar/9244773515.html}{TF5430 TwinCAT 3 Planar Motion --- Introduction}
  \item \href{https://infosys.beckhoff.com/content/1033/tc3_installation/6162675979.html}{Installing TwinCAT 3 Engineering and Runtime}
  \item \href{https://infosys.beckhoff.com/content/1033/tc3_userinterface/2867291659.html}{Information bar (stan systemu TwinCAT)}
  \item \href{https://infosys.beckhoff.com/content/1033/tc3_system/2525041803.html}{Your first TwinCAT 3 PLC project}
  \item \href{https://infosys.beckhoff.com/content/1033/tc3_system/5211863691.html}{Core Management (Real-Time: rdzenie, Base Time)}
\end{itemize}

\chapter{Konfiguracja oraz włączenie systemu}
\textit{Ćwiczenia: \textbf{inicjalizacja systemu}, \textbf{lewitacja}}

\subsection*{Backup}
\begin{itemize}
  \item \textbf{STWORZENIE BACKUP} --- student musi umieć tworzyć backup.
  \item Trzeba nauczyć dobrych praktyk: zawsze przed wejściem do środowiska produkcyjnego wykonujemy backup.
\end{itemize}

\textbf{Backup systemu na sterowniku Beckhoff (TwinCAT 3, Windows) --- skrótowo}\par
Backup składa się z trzech części; zawsze robimy je razem:
\begin{enumerate}
  \item \textbf{Archiwum projektu (źródła)} --- w drzewie projektu prawy klik na projekcie PLC $\rightarrow$ \texttt{Save <nazwa projektu> as Archive...}. Wynikiem jest plik archiwum (u nas \texttt{.tnzip}), z którego projekt można odtworzyć na innym komputerze.
  \item \textbf{Boot project} --- projekt na sterowniku musi odpowiadać projektowi w środowisku. Zapisz zmiany (czerwona dyskietka przy obiektach oznacza niezapisane zmiany), a następnie na węźle PLC wybierz \texttt{Activate Boot Project}, zaloguj się i uruchom PLC.
  \item \textbf{Kopia folderu Boot ze sterownika} (przez RDP lub udział sieciowy): w TwinCAT 3.1.4026 i nowszych \texttt{C:\textbackslash ProgramData\textbackslash Beckhoff\textbackslash TwinCAT\textbackslash 3.1\textbackslash Boot\textbackslash}. Folder \texttt{Plc} zawiera dane boot aplikacji PLC, a opcjonalnie \texttt{CurrentConfig} --- źródła projektu, jeśli są przechowywane na sterowniku.
\end{enumerate}

\uwaga{Zmienne \texttt{PERSISTENT} zapisują się w pliku \texttt{.bootdata} w folderze Boot (poprzednią wersję system zmienia na \texttt{.bootdata-old}). Aby mieć w backupie aktualne wartości, wywołaj przed kopiowaniem \texttt{FB\_WritePersistentData} z biblioteki \texttt{Tc2\_Utilities}.}

\textbf{Przywracanie:} rozpakuj archiwum projektu na komputerze, zawartość skopiowanego folderu \texttt{Plc} wstaw do folderu \texttt{Boot\textbackslash Plc} na sterowniku (zastępuje istniejącą), po czym zrestartuj TwinCAT.

\uwaga{Osobny produkt \textbf{TwinCAT Backup} (usługa \texttt{TcBackupServer}, archiwa \texttt{*.tbs}, profile z harmonogramem) wg Infosys należy do dokumentacji \textbf{TS1150 | TwinCAT 2 Backup}, więc dla TwinCAT 3 nie jest tu opisany. Źródła: Infosys Beckhoff, rozdziały \emph{Performing a PLC update}, \emph{Transferring a PLC project}, \emph{PLC Project Settings}.}
\subsection*{Pakiety i biblioteki}
\begin{itemize}
  \item Sprawdzenie i pokazanie działania Package Managera po stronie Beckhoffa.
  \item Przedstawienie bibliotek.
\end{itemize}
\subsection*{Projekt i sprzęt}
\begin{itemize}
  \item Utworzenie projektu, zdefiniowanie obszaru roboczego XPlanar.
  \item Student musi wiedzieć, jak wygląda tile.
  \item Wytłumaczenie po krótce, jak działa EtherCAT --- \textbf{max 6 tile'i na jeden EtherCAT}.
  \item Wsadzenie moverów, odczytanie BTN.
\end{itemize}
\subsection*{Odpalenie moverów}
\begin{itemize}
  \item Movery można odpalić tylko z przechyleniem do 10 stopni w osi C.
  \item Wspomnieć o ograniczeniu: nie znamy dokładnie, w jakiej orientacji jest mover \textit{(odnośnik do strony Beckhoffa; w dalszej części kursu nie będzie to wspominane)}.
  \item \textbf{Z moverUI:}
  \begin{itemize}
    \item pokazanie, jak działa JOG,
    \item pokazanie punktów rotacji, zmiany Z i innych możliwości movera \textit{(do sprawdzenia: czy na dużych moverach są punkty rotacji)},
    \item pokazanie ruchu w osi Z \textit{(można też spróbować A i B, ale nie wiem po co)}.
  \end{itemize}
  \item Na końcu: pokazanie, że movery można enable'ować przez kod --- wytłumaczenie cyklu.
\end{itemize}
\textbf{Zadanie} \textit{(propozycja)}: uruchom system w symulacji (\texttt{BasicSimulation}), zaloguj się do PLC, włącz movery z kodu (\texttt{EnableMovers}), unieś jeden mover w Z na 5 mm i opuść go z powrotem; podglądaj pozycję w \texttt{P\_ActPosition} i w Configuratorze (Live View).

\begin{itemize}
  \item \textbf{Lewitacja:} funkcje do kontroli w osi Z (\texttt{MoveZ}, zakres w konfiguracji osi: od -1 do 10 mm) \textit{(przeniesione z zagadnienia 3)}.
  \item Metoda do usuwania i dodawania moverów + wytłumaczenie, po co jest potrzebna \textit{(przeniesione z zagadnienia 3)}.
  \item Pokazanie, jak to działa dla \textbf{2 wariantów bibliotek}:
  \begin{itemize}
    \item \texttt{Tc3\_XPlanarUtility} --- ogólnie dostępna,
    \item \texttt{Tc3\_XPlanarStandard} --- poza repozytorium Beckhoffa, udostępniana na prośbę.
  \end{itemize}
\end{itemize}
\subsection*{Materiały do przekazania}
\begin{itemize}
  \item Do tego etapu: GitHub, nagranie oraz docs w PDF.
\end{itemize}

\subsection*{Dokumentacja Beckhoff (Infosys)}
\begin{itemize}
  \item \href{https://infosys.beckhoff.com/content/1033/tc3_installation/6162675979.html}{Installing TwinCAT 3 Engineering and Runtime}
  \item \href{https://infosys.beckhoff.com/content/1033/tc3_userinterface/2867291659.html}{Information bar (stan systemu TwinCAT)}
  \item \href{https://infosys.beckhoff.com/content/1033/tc3_plc_intro/2953964811.html}{Command Activate configuration}
  \item \href{https://infosys.beckhoff.com/content/1033/xts_software/10484022667.html}{Activating simulation mode}
  \item \href{https://infosys.beckhoff.com/content/1033/xts_software/11288921739.html}{Parameter (Init) --- OperationMode}
  \item \href{https://infosys.beckhoff.com/content/1033/tf5890_tc3_xplanar/19457544715.html}{TwinCAT 3 XPlanar Utility}
\end{itemize}

\chapter{Ruch wysterowany}
\textit{Ćwiczenie: \textbf{ruch w 6 stopniach swobody (6 DoF)}}

\subsection*{Start}
\begin{itemize}
  \item Z GitHuba, z brancha 3, pobrać działający projekt albo kontynuować kolejny.
\end{itemize}
\subsection*{Cel dydaktyczny}
\begin{itemize}
  \item Student powinien rozumieć, \textbf{dlaczego mamy cykl}.
  \item Pokazanie studentowi, jak ruszyć jednym moverem.
  \item Student musi umieć dodać w pętli \texttt{FOR} kolejne movery --- zrozumieć, że tak się je deklaruje i że są aktualizowane z cyklu.
\end{itemize}
\subsection*{Ruch zadany w 6 DoF --- free move}
Każda oś ma osobną komendę na obiekcie movera (\texttt{FB\_XPlanarMoverStandard}):

\begin{center}\small\begin{tabular}{@{}lll@{}}\toprule
Oś & Komenda & Uwagi \\
\midrule
X, Y (+C) & \texttt{MoveToPosition} & ruch po płaszczyźnie; zadaje X/Y/C naraz \\
C (obrót) & \texttt{MoveC} & tylko w punktach obrotu (co 120 mm); \texttt{P\_AtRotationPoint} \\
Z (wysokość) & \texttt{MoveZ} & zakres od -1 do 10 mm \\
A (przechył wokół X) & \texttt{MoveA} & zakres $\pm$6$^\circ$ \\
B (przechył wokół Y) & \texttt{MoveB} & zakres $\pm$6$^\circ$ \\
\bottomrule\end{tabular}\end{center}
\begin{itemize}
  \item **A i B nie są w \texttt{MoveToPosition}** --- zadaje się je osobnymi komendami \texttt{MoveA} / \texttt{MoveB} (parametry: pozycja docelowa w stopniach, prędkość, przyspieszenie, opóźnienie, jerk), tak samo jak \texttt{MoveZ}.
  \item Zastosowanie przechyłu: np. kompensacja przechyłu przy płynach, ustawienie narzędzia pod kątem.
\end{itemize}
\subsection*{Zadanie}
Wykonaj ruch 4 moverów w tył i w przód w osi Y.

\begin{itemize}
  \item Punkty startowe: \texttt{Y = 360}; \texttt{X} pierwszego = \texttt{120}, ostatniego = \texttt{840}.
  \item Po tym etapie student powinien rozumieć działanie cyklu, jak dodać movery itd.
  \item \textit{(Rozszerzenie 6 DoF --- do decyzji: w punktach końcowych podnieść movery w Z i przechylić w A/B.)}
\end{itemize}

\subsection*{Dokumentacja Beckhoff (Infosys)}
\begin{itemize}
  \item \href{https://infosys.beckhoff.com/content/1033/tf5430_tc3_xplanar/9335871883.html}{Limits and options of the motion commands}
  \item \href{https://infosys.beckhoff.com/content/1033/tf5430_tc3_xplanar/13424733067.html}{Example: Activating a Planar part position and moving a Planar mover}
  \item \href{https://infosys.beckhoff.com/content/1033/tf5430_tc3_xplanar/9334645515.html}{Example: Creating and moving a Planar mover with External Setpoint Generation}
\end{itemize}

\chapter{Tracki i omijanie kolizji}
\textit{Ćwiczenie: \textbf{omijanie kolizji} (tracki jako podstawa)}

\subsection*{Tracki}
\textbf{Start}

\begin{itemize}
  \item Pobieramy z GitHuba \textit{(do sprawdzenia: czy da się to dodać)} wersję kodu, w której jest dodany cały blok Motion, Mover Preparation itd.
  \item Student powinien się zaznajomić, co dokładnie dają te bloki.
\end{itemize}
\textbf{Koncept tracków}

\begin{itemize}
  \item Pokazanie konceptu tracków poprzez dokumentację --- jak go dodać.
  \item Przedstawienie ograniczeń (produkt nie jest jeszcze w 100\% dopracowany).
  \item Skupienie się na razie na \textbf{liniach}.
  \item Zlinkowanie tracków.
  \item \textbf{Group} --- deklarujemy już tutaj, bo tracki zawsze dodaje się do grupy (\texttt{EnableGroup} $\rightarrow$ \texttt{AddEnvironmentToGroup} $\rightarrow$ \texttt{EnableTracks} $\rightarrow$ \texttt{AddTracksToGroup} $\rightarrow$ \texttt{EnableMovers} $\rightarrow$ \texttt{AddMoversToGroup}). Szczegóły działania grupy --- poniżej, w części o kolizjach.
\end{itemize}
\textbf{Pokaz w UI}

\begin{itemize}
  \item Jak wejść na track.
  \item Dlaczego \textbf{keep orientation} jest ważne.
  \item JOG po tracku.
  \item Metoda \textbf{recovery}, która pozwala dołączyć do tracku \textit{(powiedzieć o ograniczeniach --- czy jest w stanie bezproblemowo nie uderzyć w inne movery)}.
\end{itemize}
\subsection*{Omijanie kolizji --- Group}
\begin{itemize}
  \item \textbf{Group} --- jak działa; odesłanie do linku Beckhoffa.
  \item Environment jako granica obszaru ruchu.
  \item Pokazanie na zasadzie kolejki \textit{(do sprawdzenia: czy nie trzeba kolejnego brancha na GitHubie, by pokazać kolejkowanie)} --- \texttt{gap}, \texttt{P\_IsQueuing}.
  \item Błąd \textbf{33158} (Collision Avoidance odrzuciło ruch) i odczyt \texttt{stInfo} --- który obiekt zablokował.
  \item Strefa wspólna --- materiał z rozdziału 08 pracy.
\end{itemize}
\subsection*{Zadania}
\textbf{Zadanie 1 --- jeden track, dwa movery:}

\begin{itemize}
  \item Pokazanie, że na tracka można wejść \textbf{ręcznie} (z UI).
  \item Zapisanie kodu dla jednego tracka (budowa geometrii, linkowanie, dodanie do Group), dołączenie do niego 2 moverów i wysterowanie jazdy w kółko --- jak w docsach Beckhoffa.
\end{itemize}
\textbf{Zadanie 2 --- dwa tracki z punktem przecięcia:}

\begin{itemize}
  \item 2 movery jadące na \textbf{oddzielnych trackach}, które przechodzą przez jeden punkt przecięcia (łatwe do zrobienia).
  \item Obserwacja zatrzymania, kolejkowania i wznowienia przy kolizji w punkcie przecięcia.
\end{itemize}

\subsection*{Dokumentacja Beckhoff (Infosys)}
\begin{itemize}
  \item \href{https://infosys.beckhoff.com/content/1033/tf5430_tc3_xplanar/9247679243.html}{Planar track operation modes}
  \item \href{https://infosys.beckhoff.com/content/1033/tf5430_tc3_xplanar/9021192715.html}{Planar environment}
  \item \href{https://infosys.beckhoff.com/content/1033/tf5430_tc3_xplanar/19794139531.html}{Collision avoidance between Planar movers and Planar tracks}
  \item \href{https://infosys.beckhoff.com/content/1033/tf5430_tc3_xplanar/10384430603.html}{Example: Creating and moving Planar movers with group}
  \item \href{https://infosys.beckhoff.com/content/1033/tf5430_tc3_xplanar/9253329803.html}{Example: Creating and moving Planar movers with track and group}
  \item \href{https://infosys.beckhoff.com/content/1033/tf5430_tc3_xplanar/10789602827.html}{Example: Following a Planar mover through a Track Network}
\end{itemize}

\chapter{Dynamiczne trasowanie moverów --- stacje}
\textit{Ćwiczenie: \textbf{dynamiczne trasowanie moverów} (najtrudniejsze)}

\begin{itemize}
  \item Stacje jako punkty decyzyjne --- mover po obsłudze dostaje kolejny cel w zależności od stanu systemu.
  \item Wzorzec ze starter kitu: \texttt{FB\_Station\_01..03} (stacje procesowe, wybór następnej stacji), \texttt{FB\_Station\_04\_SW\_Separate} (zwrotnica 1$\rightarrow$2), \texttt{FB\_Station\_05\_SW\_Merge} (zwrotnica 2$\rightarrow$1 z priorytetem).
  \item Listy kolejek stacji (\texttt{FB\_List}), \texttt{P\_ProcessRequest} / \texttt{P\_ProcessClearance}.
  \item \texttt{FB\_MoverTarget} --- przydział najbliższego movera do pozycji.
\end{itemize}
\textbf{Zadanie} \textit{(propozycja)}: na torach z zagadnienia 3 dodać 2 stacje i zwrotnicę; movery cyklicznie trafiają do wolnej stacji.


\subsection*{Dokumentacja Beckhoff (Infosys)}
\begin{itemize}
  \item \href{https://infosys.beckhoff.com/content/1033/tf5430_tc3_xplanar/9244773515.html}{TF5430 TwinCAT 3 Planar Motion --- Introduction}
  \item \href{https://infosys.beckhoff.com/content/1033/tf5430_tc3_xplanar/10789602827.html}{Example: Following a Planar mover through a Track Network}
\end{itemize}

\end{document}
